\chapter{Épreuve chronométrée 10}\label{ch:iv1:timed-test-10}

\begin{ivtrains}
La réaliser d'un seul tenant, sans notes ni calculatrice.
\end{ivtrains}

\ivtestsection{A}{6}{40}

\begin{question}\label{iq:iv1:timed-test-10:1}
$X$, $Y$ et $Z$ sont indépendantes et uniformes sur $[0, 1]$. Quelle est la probabilité que $X + Y + Z$ se situe entre 1 et 2~?
\end{question}

\begin{question}\label{iq:iv1:timed-test-10:2}
Cinq balles différentes sont lancées indépendamment et au hasard dans 3 urnes. Quelle est la probabilité qu'aucune urne ne soit vide~?
\end{question}

\begin{question}\label{iq:iv1:timed-test-10:3}
Un rendement a une moyenne de 10 et une variance de 4. Majorez la probabilité d'un écart d'au moins 5 par rapport à la moyenne à l'aide de l'inégalité de Bienaymé-Tchebychev. Quelle est cette probabilité si le rendement est normal, sachant que $P(Z > 2.5) = 0.0062$~?
\end{question}

\begin{question}\label{iq:iv1:timed-test-10:4}
Trois dés ont les faces suivantes~: A porte (2, 2, 4, 4, 9, 9), B porte (1, 1, 6, 6, 8, 8) et C porte (3, 3, 5, 5, 7, 7). Chacun a le même total. Quelles sont les probabilités que A batte B, que B batte C et que C batte A, en un seul lancer chacun contre l'autre~?
\end{question}

\begin{question}\label{iq:iv1:timed-test-10:5}
Deux lignes de bus arrivent selon des processus de Poisson indépendants, avec des intervalles moyens de 10 et de 20 minutes. Vous arrivez à un instant aléatoire et prenez le premier bus. Quelle est la probabilité de prendre la première ligne~? Quelle est votre attente moyenne~?
\end{question}

\begin{question}\label{iq:iv1:timed-test-10:6}
Une année sur une autre planète compte 20 jours, et les anniversaires sont indépendants et uniformes. Quelle est la plus petite taille de groupe pour laquelle la probabilité d'un anniversaire commun dépasse $\frac12$~?
\end{question}

\ivtestsection{B}{6}{60}

\begin{question}\label{iq:iv1:timed-test-10:7}
$X$ est uniforme sur $[0, 1]$. Sachant $X$, $Y$ est uniforme sur $[0, X]$. Que vaut $\mathbb{E}[X \mid Y = \frac12]$~? Vous pouvez utiliser $\ln 2 = 0.693$.
\end{question}

\begin{question}\label{iq:iv1:timed-test-10:8}
Cinq dés équilibrés sont lancés. Quelle est la probabilité que leur total soit d'au moins 18~? Avec six dés équilibrés, quelle est la probabilité d'un total d'au moins 21, sachant que $4\,332$ des $46\,656$ issues totalisent exactement 21~?
\end{question}

\begin{question}\label{iq:iv1:timed-test-10:9}
Un flux contient un nombre inconnu d'éléments. Vous conservez un élément en mémoire~: le premier au départ, puis vous le remplacez par le $i$-ième élément avec la probabilité $\frac1i$. Quelle est la probabilité que le troisième des 5 éléments soit celui conservé à la fin~? Comment modifier cette méthode pour en conserver 2, et quelle est alors la probabilité pour chacun des éléments~?
\end{question}

\begin{question}\label{iq:iv1:timed-test-10:10}
Une main de cinq cartes est tirée d'un jeu standard de 52 cartes, qui compte $2\,598\,960$ mains possibles. Quelle est la probabilité d'obtenir exactement un brelan, c'est-à-dire trois cartes d'un même rang et deux autres cartes de deux rangs différents entre eux et différents du premier~?
\end{question}

\begin{question}\label{iq:iv1:timed-test-10:11}
Sur 100 lancers d'une pièce équilibrée, quelle est la probabilité que le nombre de piles diffère de 50 d'au moins 10~? Utilisez l'approximation normale avec correction de continuité, sachant que $P(Z > 1.9) = 0.0287$.
\end{question}

\begin{question}\label{iq:iv1:timed-test-10:12}
Un action vaut 100. Un call européen de strike 105 et d'échéance un an coûte 10, et le taux d'intérêt est nul. Que doit valoir un put européen de même strike et de même échéance~? S'il se cote à 14, quelle transaction garantit un profit, et de combien~?
\end{question}
